% @econ-contract section
% id: RESULTS_HORIZONS
% role: Establish on the class what foresight adds to the myopic benchmark (thm:regime-separation: defiance, a sanction above the myopic sanction, a state path above the benchmark's, a positive steady state, also at every full-defiance equilibrium); split the difference into the Target's part (myopic relief against the entrenchment condition) and the Sender's part (the foresight premium and the fixed-policy comparison); then lock-in, the role of the Target's patience, and the minmax bound and uniqueness in words. The cases with one government forward-looking are in the discussion (sec:discussion:horizons).
% question_answered: What does foresight add to the myopic benchmark, what does each government's incentive contribute, and which conclusions extend beyond the constructed equilibria?
% required_inputs: MODEL_ENVIRONMENT (objects, the class, the myopic benchmark and its assumptions, the equilibrium concept).
% output_handoff: The comparison of the forward-looking equilibrium with the benchmark and its quantified scope, used by the welfare section (the two constructed equilibria) and the discussion (the four cases, scope and persistence).
% claim_ceiling: The constructed forward-looking equilibrium against the myopic benchmark, and the same conclusions at every full-defiance equilibrium; a same-policy foresight premium; lock-in under its stated conditions; existence inheritance under monotonicity and balance; an exact threshold only on the tight collapse class; a value bound in every top-selector equilibrium and identification only in the nondecreasing-rule class (stated in words).
% exclusions: The proofs (Online Appendix, one footnote per subsection); the four cases (sec:discussion:horizons); the component statements, the strict comparison, the minmax and uniqueness propositions and the relief/defiance inequalities (Online Appendix, sec:moved-main-2026-10-13); welfare.
% exact_ref: tex-label:sec:results
% support_status: VERIFIED
% status: HUMAN_CONFIRMED
% human_confirmed_at: 2026-10-01
% decision_ref: suggestions/2026-10-01-manuscript-rebuild/01_STRUCTURE_PROPOSAL.md#author-decisions-recorded
% confirmation_scope: Section role, claim ceiling, and the four short propositions (D2, D5).
% reopen_if: A principal result, its hypotheses, or the order of the moves changes.
% amended_at: 2026-10-11 (lock-in subsection and prop:lock-in; decision_ref: suggestions/2026-10-10-jet-submission/DECISIONS.md#lock-in-in-the-main-text); 2026-10-11 (no 'horizon'; no wording about what the paper is not; decision_ref: suggestions/2026-10-10-jet-submission/DECISIONS.md#plain-framing-2026-10-11); 2026-10-13 (exposition: results in order of importance, Tier-3 statements moved to the Online Appendix, prop:two-sided-one-equilibrium moved in, one proof footnote per subsection; decision_ref: suggestions/2026-10-10-jet-submission/DECISIONS.md#exposition-2026-10-13); 2026-10-15 (the dual-basin economy and prop:two-sided-one-equilibrium moved back to the Online Appendix, app:dual-basin-economy, on the author's word; decision_ref: suggestions/2026-10-10-jet-submission/DECISIONS.md); 2026-10-16 (layout decisions D1-D5: the main comparison is the forward-looking equilibrium against the myopic benchmark of the model section, the two governments' parts follow it, and the four cases moved to the discussion; decision_ref: suggestions/2026-10-10-jet-submission/DECISIONS.md#layout-myopic-benchmark-first-2026-10-16); 2026-10-16 (structure plan S1-S7 against the cited papers; decision_ref: suggestions/2026-10-10-jet-submission/DECISIONS.md#structure-plan-s1-s7-2026-10-16); 2026-10-19 (coauthor wording review: wording of the coauthor's line-by-line revision applied where not wrong; sender/target in lower case; non-compliance for defiance; decision_ref: suggestions/2026-10-10-jet-submission/DECISIONS.md#coauthor-wording-review-2026-10-19); 2026-10-20 (RESULTS_LOCKIN_READING: the absorbing basin and the lasting effect of a temporary event, Acemoglu and Robinson 2008, 2001; decision_ref: suggestions/2026-10-10-jet-submission/DECISIONS.md#acemoglu-comparison-and-model-exposition-2026-10-20); 2026-10-20 (subsection openings in the style of Acemoglu and coauthors, N187; the claims unchanged; decision_ref: suggestions/2026-10-10-jet-submission/DECISIONS.md#acemoglu-style-subsection-openings-2026-10-20); 2026-10-20 (plain-language reminders at the mathematics, N188; formal statements unchanged; decision_ref: suggestions/2026-10-10-jet-submission/DECISIONS.md#plain-language-reminders-at-the-mathematics-2026-10-20); 2026-10-20 (equation prose in the forms of Acemoglu and Robinson 2008, N189; decision_ref: suggestions/2026-10-10-jet-submission/DECISIONS.md#equation-prose-in-the-style-of-acemoglu-and-robinson-2026-10-20); 2026-10-20 (two footnote pointers follow the low-basin theorem and post-relief persistence into the persistence section of the Online Appendix; decision_ref: suggestions/2026-10-10-jet-submission/DECISIONS.md#appendix-needed-versus-derivative-2026-10-20)
% @end-econ-contract
\section{Equilibrium analysis: foresight and sanction persistence}
\label{sec:results}
\label{sec:existence}

% @econ-contract paragraph
% id: RESULTS_ROADMAP
% role: Say what the section establishes, and keep existence, uniqueness and value bounds apart.
% depends_on: RESULTS_HORIZONS
% handoff: The main comparison.
% support_status: NOT_APPLICABLE
% status: PROPOSED
% @end-econ-contract
\central{This section compares forward-looking governments with the myopic
benchmark.} \Cref{sec:existence:placable} shows that foresight turns relief into
non-compliance and keeps sanctions and the political state positive.
\Cref{sec:results:wedge} bounds the sender's part of that difference. \Cref{sec:results:lockin} gives conditions under which a high
political state stays high, \Cref{sec:results:patience} asks how patient the
target must be and collects two further consequences
of non-compliance. Each subsection opens with a footnote that says where its proofs are.
Existence, uniqueness within a stated set and value bounds are kept separate, each
under the hypotheses of its own result.

\subsection{Comparison with the myopic benchmark}
\label{sec:existence:placable}
\label{sec:results:main}

% @econ-contract paragraph
% id: RESULTS_CLASS_PURPOSE
% role: Explain how the full-defiance equilibrium is constructed (prop:placable-family, stated in the Online Appendix), and state thm:regime-separation, the comparison of the forward-looking equilibrium with the myopic benchmark (the main theorem since 2026-10-16).
% depends_on: RESULTS_ROADMAP
% handoff: The table and the reading.
% exact_ref: tex-label:thm:regime-separation
% exact_ref: tex-label:prop:placable-family
% exact_ref: tex-label:sec:moved-main-2026-10-13
% exact_ref: tex-label:app:proofs:placable
% support_status: VERIFIED
% status: PROPOSED
% @end-econ-contract
\central{We now turn to the main result: when both governments are forward-looking, the same
economy supports an equilibrium with persistent non-compliance and indefinite sanctions.}\footnote{\Cref{thm:regime-separation}
is proved in panel~F of \Cref{app:proofs:placable} of the Online Appendix, and the
full-non-compliance equilibrium it uses (\Cref{prop:placable-family}) is constructed in
the same appendix; the component statements are in
\Cref{sec:moved-main-2026-10-13}.} Stage
payoffs, feasible sets, and the transition law are identical to the benchmark; only
the discount factors are positive. Superscripts index the case: $00$ for the myopic
benchmark, $11$ for the case with two forward-looking governments. The non-compliance
equilibrium is constructed in two steps. First, the sender's value under full
non-compliance is derived as the fixed point of its Bellman operator; at every state,
the forward-looking sanction is then strictly higher than the corresponding myopic
sanction $\sigM(x,0)$, the sanction a myopic sender sets at state $x$ against full
non-compliance. Second, under the entrenchment and increasing-differences
conditions, full non-compliance is the target's unique best reply to that response, and
the two value functions paired with the full-non-compliance policy constitute an
equilibrium in $\mathcal E$. Recall that $\mathcal E$ is the set of stationary feedback
Stackelberg equilibria (\Cref{sec:model:equilibrium}) and that $\underline s$ is the least
sanction a myopic sender imposes on a non-compliant target (\Cref{sec:model:objectives}).

\begin{theorem}[Foresight vs.\ the myopic benchmark]
\label{thm:regime-separation}
For every economy in the class satisfying
Assumptions~\ref{ass:stage-interiority}--\ref{ass:dynamic-complementarity}, there
exists an equilibrium in $\mathcal E$ such that:
\begin{enumerate}[label=(\roman*)]
\item the target maintains full non-compliance at every state, and the sender's
sanction level lies strictly above the myopic sanction $\sigM(x,0)$;
\item from every initial state the amount by which the political state falls short of
$\underline s$ shrinks at least geometrically; starting from a common initial state,
the state lies strictly above the path of the myopic benchmark
(\Cref{prop:myopic-benchmark}) at every positive date;
\item the state converges to a steady state $x^{11}_\infty>\underline s$, at which the
sanction level equals the state level.
\end{enumerate}
\end{theorem}

% @econ-contract display
% id: TAB_MYOPIC_FORWARD
% role: Show, for the myopic benchmark and the forward-looking equilibrium, the compliance, the sanction and the long-run state.
% evidence_represented: prop:myopic-benchmark and thm:regime-separation.
% permitted_inference: The behavior and the limits of the two constructed equilibria, under the stated conditions.
% reader_takeaway: Foresight turns relief into defiance, and sanctions and the political state then stay positive.
% claim_supported: RESULTS_MAIN_READING
% source_artifact: paper/main_sections/03_equilibrium_and_comparative_statics.tex
% caption_scope: The two constructed equilibria only.
% exact_ref: tex-label:thm:regime-separation
% exact_ref: tex-label:prop:myopic-benchmark
% support_status: VERIFIED
% status: PROPOSED
% @end-econ-contract
\begin{table}[htbp]
\centering\footnotesize
\caption{The myopic benchmark and the forward-looking equilibrium of one economy:
compliance, sanction and the limit of the political state. Stage payoffs, feasible
actions and the transition are the same in both columns; the conditions of
\Cref{thm:regime-separation} are maintained.}
\label{tab:myopic-forward}
\begin{tabular}{@{}p{0.16\linewidth}p{0.38\linewidth}p{0.4\linewidth}@{}}
\toprule
 & myopic governments (\Cref{prop:myopic-benchmark}) & forward-looking governments (\Cref{thm:regime-separation}) \\
\midrule
compliance & $a^{00}(x)$, the least that ends the myopic sanction & full non-compliance (compliance $0$) \\
sanction & zero at every state & strictly above the myopic sanction $\sigM(x,0)$, and below $1$ under interiority \\
political state & falls to $0$, at least geometrically & above the benchmark's path at every positive date; converges to a steady state above $\underline s$, which may depend on the start \\
\bottomrule
\end{tabular}
\end{table}

% @econ-contract paragraph
% id: RESULTS_MAIN_READING
% role: Read the theorem against the benchmark: foresight turns relief into defiance (the investment logic of the entrenchment condition), the forward-looking Sender sanctions more than the myopic sanction, the state settles at a positive level; the conclusions hold at every full-defiance equilibrium; both discount factors move together, and where the separate roles are found.
% depends_on: RESULTS_CLASS_PURPOSE
% handoff: The two governments' parts.
% merged_from: RESULTS_FOUR_CELLS_READING (its defiance reading), 2026-10-16
% exact_ref: tex-label:thm:regime-separation
% exact_ref: tex-label:prop:myopic-benchmark
% exact_ref: tex-label:tab:myopic-forward
% support_status: VERIFIED
% status: PROPOSED
% @end-econ-contract
\central{Foresight shifts outcomes from relief to non-compliance, and under
non-compliance both sanctions and the political state remain positive.} Conclusion (ii)
implies that a state that starts below $\underline s$ rises toward $\underline s$ or above
it, because its shortfall from $\underline s$ shrinks at least geometrically. Conclusion
(iii) says that $x^{11}_\infty$, the long-run political state when both governments look
ahead, lies above $\underline s$; at this state the sanction equals the state, so by
\textup{(F4)} the state stays there. In the myopic
benchmark of \Cref{prop:myopic-benchmark}, the target accommodates and the state
erodes to zero (\Cref{tab:myopic-forward}). The forward-looking target instead chooses
non-compliance because, under the entrenchment condition, the entrenchment generated
by sanctions is worth at least as much as the immediate pain they impose, making it
unprofitable to trade compliance for relief. For its part, the forward-looking sender
sanctions a non-compliant target more heavily than a myopic sender would. The state
then converges to a positive level from any initial condition, declining toward it
from sufficiently high starting points. Under interiority (\Cref{ass:interiority}), the
sender never imposes full sanctions. Conclusions (i)--(iii) hold at every equilibrium
in $\mathcal E$ with full target non-compliance---not only at the constructed
equilibrium---because all such equilibria share the same constructed values and
sanction rule. The steady-state limit may depend on initial conditions unless the
equilibrium law of motion is contractive in state distances
(\Cref{sec:results:lockin}).

\central{The theorem changes both discount factors at once, holding every other
primitive fixed.} It shows that foresight makes sanctions persist; it does not by
itself say which government's foresight does what. \Cref{sec:results:wedge} bounds
the sender's part of the difference, and
\Cref{sec:discussion:horizons} solves the cases in which only one government looks
forward.

\subsection{The sender's foresight premium}
\label{sec:results:wedge}
\label{sec:compstat}

% @econ-contract paragraph
% id: RESULTS_WEDGE
% role: Give the Sender's part: the closed-form lower bound on the foresight premium, what each term measures, and the fixed-policy comparison behind it (the foresight premium along paths and the long-run gap over the Target-only case moved to sec:discussion:horizons, 2026-10-16).
% depends_on: RESULTS_MAIN_READING
% handoff: The fixed-policy comparison.
% exact_ref: tex-label:eq:sender-wedge
% exact_ref: tex-label:rem:wedge-sharp
% support_status: VERIFIED
% status: PROPOSED
% @end-econ-contract
\central{The comparison so far concerns the equilibrium as a whole; we next isolate the sender's
contribution, the foresight premium: the amount by which the sender's forward-looking
sanction exceeds its myopic sanction at a given compliance level.} This premium has a closed-form lower bound at
full non-compliance that increases with the sender's stake in the political
state.\footnote{The bound
and the attainment example are proved in \Cref{app:proofs:branch} of the Online
Appendix; the comparison at a fixed target policy and its strict form in
\Cref{app:proof-dominance}.} Under interiority, at every state the forward-looking sender's sanction
under non-compliance exceeds the myopic sanction $\sigM(x,0)$ (the myopic sender's sanction against full
non-compliance) by at least a strictly positive closed-form bound. The numerator captures the minimum increase in the
sender's discounted value from a unit of sanctions via the next-period state: the
sender's per-period stake in the state, discounted over the state's persistence rate,
multiplied by the minimum response of the next-period state to sanctions. The
denominator---the maximum curvature of the sender's sanctioning cost function---converts
this marginal gain into a sanction increase. The bound rises with the sender's
discount factor, its stake in the state, and the responsiveness of the transition to
sanctions; its general form is given by \eqref{eq:sender-wedge} in the Online Appendix. These lower bounds hold for fixed primitives. The bound is attained at one
example economy, so its coefficient cannot be uniformly scaled up across the class.
The foresight premium rests on a comparison that holds broadly outside the class. Recall that $\sigM$ is the myopic best response and
$\sigDynTop_{\Vs}$ the dynamic best response, which also counts the value $\Vs$ of the next
state; $\alpha$ is a target policy.

\begin{theorem}[Dynamic vs.\ myopic sanctions at a fixed target policy]
\label{thm:dominance}
Let $\Vs$ be a bounded continuous sender continuation value that is non-decreasing in
the political state. For any target policy $\alpha\colon[0,1]\to[0,1]$ and every state
$x\in[0,1]$,
\begin{equation}
\sigM(x,\alpha(x))
\le
\sigDynTop_{\Vs}(x,\alpha(x)).
\label{eq:thm-dominance-same-alpha}
\end{equation}
In particular, inequality \eqref{eq:thm-dominance-same-alpha} holds at every
equilibrium $E=(\Vt,\Vs,\alpha)\in\mathcal E$ whose sender value is non-decreasing in
the state.
\end{theorem}

% @econ-contract paragraph
% id: RESULTS_THREE_COMPARISONS
% role: Give the reason for thm:dominance, its strict form in words (prop:strict-dominance, stated in the Online Appendix), and the direct and indirect effects against the myopic equilibrium, read as the Sender's and the Target's parts of the move from the benchmark (2026-10-16).
% depends_on: RESULTS_WEDGE
% handoff: Lock-in.
% exact_ref: tex-label:thm:dominance
% exact_ref: tex-label:prop:strict-dominance
% exact_ref: tex-label:cor:pathwise-dominance
% support_status: VERIFIED
% status: PROPOSED
% @end-econ-contract
\central{The intuition relies on monotonicity: at a fixed target policy, a higher
sanction leaves a weakly higher successor state, so a non-decreasing continuation value
can only increase the largest of the optimal sanctions.} Inequality
\eqref{eq:thm-dominance-same-alpha} thus implies that a sender that values a more
entrenched adversary never sanctions less than a myopic sender facing the same
compliance. The comparison holds under general
transition rules and requires neither an equilibrium nor restriction to high states.
Every stage maximizer, not only the selected one, lies weakly below the forward-looking
sanction. The inequality becomes strict when no sanction below the greatest stage
maximizer yields a higher continuation value (automatic when the continuation value is
non-decreasing) and, for some sanction level, the discounted continuation gain from
moving to that level outweighs the stage loss relative to the greatest stage
maximizer. This version requires no concavity, differentiability, or positive
discount.

Relative to the myopic equilibrium at its own target policy, the forward-looking
equilibrium sanctions at least as heavily at every state when its sender value is
non-decreasing, its target complies no more, and the sender's myopic best response does
not increase with compliance. The difference operates through two separate channels at
each state where the forward-looking target complies no more than the myopic one: lower
target compliance can only raise the myopic best response (the indirect effect), and
sender foresight can only raise sanctions at that compliance level (the direct effect,
\Cref{thm:dominance}). At the level of the sanction rule, these two effects---the
target's part and the sender's part---sum to the total difference between
\Cref{prop:myopic-benchmark} and \Cref{thm:regime-separation}.

\subsection{Sanction lock-in and history dependence}
\label{sec:results:lockin}

% @econ-contract paragraph
% id: RESULTS_LOCKIN
% role: Say what lock-in means here, link it from the defiance result to the Sender's side, and state the lock-in proposition at its scope with its own symbols defined.
% depends_on: RESULTS_THREE_COMPARISONS
% handoff: The conditions one by one.
% exact_ref: tex-label:thm:persistence
% exact_ref: tex-label:thm:shock
% exact_ref: tex-label:sec:assumptions
% support_status: VERIFIED
% status: PROPOSED
% @end-econ-contract
\central{The results so far describe constructed equilibria. We now show that, under the
conditions stated below, persistence sustains itself once it starts, a property we call
lock-in: once the political state reaches a high level, equilibrium play keeps it there,
the state never exits a high basin, and sanctions never decline.}\footnote{\Cref{prop:lock-in} restates the persistence
and temporary-shock theorems of \Cref{app:persistence} in the Online Appendix,
which also gives their value sets, the constants of the level condition and the
shock length; the low-basin counterpart is
\Cref{thm:lowbasin}, and an equilibrium with both basins in
\Cref{app:dual-basin-economy}.} Whereas the non-compliance result focuses on the target
side, lock-in incorporates the sender's incentives. Under the conditions below,
raising sanctions to match the state is profitable for the sender across the entire high
basin---from a threshold $\bar x$ up to the upper bound---and a sanction at or above the
current state prevents the state from falling. The result requires none of the
class-wide bounds: it holds at every equilibrium of any economy satisfying its
conditions. A temporary shock---a finite period of exogenously set sanctions---can move
the state into the high basin. In the proposition, $E=(\Vt,\Vs,\alpha)$ is an equilibrium:
the two governments' value functions and the target's policy.

\begin{proposition}[Lock-in]
\label{prop:lock-in}
Let $E=(\Vt,\Vs,\alpha)\in\mathcal E$ and let $[\bar x,1]$ be a high basin of states.
Suppose that the equilibrium value functions lie within closed sets of bounded
continuous functions, with sender values Lipschitz-continuous in the state, and that
the target policy is the one induced by these value functions. Suppose further that, at
every pair of values in those sets (with the induced target policy substituted):
\begin{enumerate}[label=(\alph*)]
\item \emph{monotonicity}: on the high basin, the sender's objective satisfies the
interval dominance order in the sanction and the state \citep{QuahStrulovici2009}, so the largest of its optimal sanctions rises with the state; and
\item \emph{level condition}: at every state in the basin, raising sanctions from any
lower level up to the current state level increases the sender's current payoff by more
than the maximum possible change in its discounted continuation value.
\end{enumerate}
Then from every initial state $x_0\in[\bar x,1]$, the equilibrium path
$s_t=\sigDynTop_{\Vs}(x_t,\alpha(x_t))$ and $x_{t+1}=\F(x_t,s_t)$ satisfies for all
$t\ge0$:
\begin{enumerate}[label=(\roman*)]
\item $x_t\in[\bar x,1]$;
\item $s_t\ge\bar x$;
\item $s_t\le s_{t+1}$.
\end{enumerate}
If in addition the transition satisfies \textup{(F1)--(F4)} and full sanctions raise
every state below the basin, then from any initial state a temporary shock holding
sanctions at the maximum level for a sufficient duration moves the state into
$[\bar x,1]$ once the shock ends, and properties \textup{(i)--(iii)} hold along the
equilibrium path of $E$ restarted from that state.
\end{proposition}

% @econ-contract paragraph
% id: RESULTS_LOCKIN_READING
% role: Read the two parts of prop:lock-in against the political economy of persistent power: the high basin as an absorbing region (Acemoglu and Robinson 2008) and a temporary event with a lasting effect (Acemoglu and Robinson 2001) (N184, N186, 2026-10-20).
% depends_on: RESULTS_LOCKIN
% handoff: The conditions of the proposition.
% exact_ref: bib:AcemogluRobinson2008#p=279
% source_role: bib:AcemogluRobinson2008#p=279 => Corollary 2: when investment in de facto power is not always profitable for the elite, 'democracy can become an absorbing state' (read 2026-10-20, N184); comparison, not support.
% exact_ref: bib:AcemogluRobinson2001#p=939
% source_role: bib:AcemogluRobinson2001#p=939 => Because 'the threat of revolution is often only transitory', the elite is 'forced to make a credible commitment to future income redistribution', which extending voting rights achieves (read 2026-10-20, N186); comparison, not support.
% support_status: VERIFIED
% status: PROPOSED
% @end-econ-contract
\central{Property~(i) makes the high basin an absorbing region of the equilibrium
path.} In words, (i) the state never falls below $\bar x$; (ii) the sanction never falls
below $\bar x$; (iii) the sanction never falls from one date to the next. Democracy can become an absorbing state in the same sense in
\citet[Corollary~2, p.~279]{AcemogluRobinson2008}, when investment in de facto power is
no longer always profitable for the elite. The entry result gives a temporary event a
lasting effect, as in \citet[p.~939]{AcemogluRobinson2001}, where a transitory threat
of revolution forces the elite to extend the franchise; there the lasting change runs
through the credibility of concessions, here through the political state that the
shock moves.

% @econ-contract paragraph
% id: RESULTS_LOCKIN_CONDITIONS
% role: Discuss the conditions of prop:lock-in one by one, say which conclusion uses which, and state the weaker form that holds on the class at every full-defiance equilibrium.
% depends_on: RESULTS_LOCKIN
% handoff: Low sanctions entrenched, and persistence without hysteresis.
% exact_ref: tex-label:thm:persistence
% exact_ref: tex-label:thm:finite-episode
% support_status: VERIFIED
% status: PROPOSED
% @end-econ-contract
\central{The conclusions of \Cref{prop:lock-in} rely on different subsets of its
conditions.} For economies in the class satisfying interiority and the lifting
condition, a weaker form of lock-in holds at every full-non-compliance equilibrium. The
set-based conditions restrict equilibrium value functions to sets of bounded functions
with Lipschitz sender values, and require the target policy to be the one induced by
the value pair. The monotonicity condition ensures the sender's response (with the
target policy substituted) is non-decreasing in the state across the basin. The level
condition requires that, at each basin state, raising sanctions from any lower level up
to the state increases the sender's current payoff by more than any associated change in discounted continuation value. Both conditions are imposed pointwise across all
value pairs in the sets. The no-exit properties (i) and (ii) rely only on the set
conditions and the level condition; the non-decreasing sanction property (iii)
additionally requires the monotonicity condition. No assumption on how the target ranks
sanction levels is needed. The entry result additionally requires conditions \textup{(F1)--(F4)}
and that full sanctions raise every state below $\bar x$; the Online Appendix provides
an explicit sufficient shock duration.

For the class of economies satisfying interiority and the lifting condition, at every
equilibrium with full target non-compliance, once a temporary shock moves the state to
or above $\underline s$ (the least sanction a myopic sender imposes on a non-compliant target), the state remains there and
sanctions never fall below $\underline s$.

% @econ-contract paragraph
% id: RESULTS_LOCKIN_PEACE
% role: State the low-basin dual and initial-condition selection in words, and separate persistence from dependence on initial conditions.
% depends_on: RESULTS_LOCKIN_CONDITIONS
% handoff: Target patience and the existence of defiance.
% exact_ref: tex-label:thm:lowbasin
% exact_ref: tex-label:cor:initial-condition-selection
% exact_ref: tex-label:prop:activation-not-hysteresis
% support_status: VERIFIED
% status: PROPOSED
% @end-econ-contract
\central{The argument also runs in reverse: the political state can entrench low
sanctions as well as high ones.} Under low-basin versions of the conditions, the state
remains in the basin and realized sanctions never rise for any starting point within
a low basin. When a single equilibrium satisfies the sender's conditions on
both a high basin and a disjoint low basin below it, the initial state selects between
the two basins within that equilibrium. Persistence can also hold independently of
initial conditions: for economies in the class satisfying interiority, the lifting
condition, and the entrenchment and increasing-differences conditions, and whose
equilibrium law of motion has an explicit Lipschitz constant strictly below one, the
constructed equilibrium has exactly one steady state. Every initial state converges to
this state, and temporary sanction shocks have no permanent effect on the long-run
limit.

% @econ-contract paragraph
% id: DISC_RELIEF_EXIT
% role: Imposed relief: the crossing, the condition on the resumed equilibrium under which de-escalation lasts, and the comparison of entry and exit durations (moved from the Discussion, 2026-10-16, structure plan S2).
% depends_on: RESULTS_LOCKIN_PEACE
% handoff: Target patience and the existence of defiance.
% exact_ref: tex-label:thm:relief-exit
% exact_ref: tex-label:prop:post-relief-persistence
% support_status: VERIFIED
% status: PROPOSED
% @end-econ-contract
\label{sec:discussion:persistence}
\central{The reverse experiment---a temporary suspension of sanctions---can reduce the
state below a given level within a computable time horizon.}\footnote{Entry and exit durations, persistence after relief, the reading of the
lock-in conditions and the example economies that meet them are in
\Cref{app:persistence,app:witness} of the Online
Appendix.} After entry into a
high basin, equilibrium play generates no endogenous exit, but relief remains
mechanically feasible: the transition law still erodes the state under low sanctions.
Whether de-escalation persists then depends on the resumed equilibrium: it persists
if the equilibrium's sanction rule remains at or below the state level for states
below that threshold. Entry and exit durations derive from different segments of the same transition
function, so they need not be equal. For any transition symmetric under reflection of
the unit square, $\F(1-x,1-s)=1-\F(x,s)$ (the state moves up toward a sanction above it in the same way as
it moves down toward one below it), such as the partial-adjustment law, the
exit-duration bound (before rounding to integer periods) exceeds the entry-duration
bound
exactly when that level lies below one half.

\subsection{Target patience and non-compliance equilibria}
\label{sec:results:patience}
\label{sec:results:bracket}

% @econ-contract paragraph
% id: RESULTS_COLLAPSE
% role: State the exact threshold discount factor for full defiance on the tight class (prop:collapse-threshold), its hypotheses in words, and what it means.
% depends_on: RESULTS_LOCKIN_PEACE
% handoff: A more patient Target keeps the full-defiance equilibrium.
% exact_ref: tex-label:prop:collapse-threshold
% exact_ref: tex-label:prop:bracket
% exact_ref: tex-label:prop:defiance-closed
% support_status: VERIFIED
% status: PROPOSED
% @end-econ-contract
\central{The analysis so far compares a myopic with a forward-looking target. A natural question
is how patient the target must be for non-compliance to prevail; the comparison extends
across discount factors and collapses to a single threshold in the linear--quadratic
special case.}\footnote{The
threshold, the general sufficient conditions for relief and for non-compliance, the
exact thresholds and the inheritance by a more patient target are proved in
\Cref{app:proofs:threshold,app:proofs:scope} of the Online Appendix, and branch
invariance in \Cref{app:proofs:branch}; the general conditions, the precise
conditions of the special case and the closed form of the threshold are stated in
\Cref{sec:moved-main-2026-10-13:bracket}.} In this special case: the transition responds to sanctions at a constant
rate, which combined with $\F(x,x)=x$ makes it affine; the curvature of the sender's
sanctioning cost is constant; the substitution rate, the target's compliance cost, its
sanction pain, and its valuation of the state are all constant; the target's payoff has
no state--sanction cross term; and the curvature bounds on the sender's value and its
state-dependent payoff are zero. The target's \emph{net cost} of sanctions equals their
pain minus the compliance cost of ending them.

\begin{proposition}[Threshold discount factor for full non-compliance]
\label{prop:collapse-threshold}
In the linear--quadratic special case, for an economy satisfying interiority, the
myopic relief condition, and the lifting condition, there exists a threshold discount
factor $\beta_T^*$ (in closed form in the Online Appendix) such that an equilibrium in
$\mathcal E$ with full non-compliance at every state exists if and only if
$\beta_T\ge\beta_T^*$. When $\beta_T^*<1$, these discount factors span the interval
$[\beta_T^*,1)$.
\end{proposition}

% @econ-contract paragraph
% id: RESULTS_THRESHOLD_READING
% role: Read the threshold in words (2026-10-16, no constants of the special case in the main text): what raises and lowers it, the case beta_T* >= 1, and why the Sender's patience does not move it.
% depends_on: RESULTS_COLLAPSE
% handoff: The more patient Target.
% merged_from: RESULTS_SENDER_BRACKET
% exact_ref: tex-label:prop:collapse-threshold
% exact_ref: bib:AcemogluJensen2015#scope=Lemma2-Prop1
% source_role: bib:AcemogluJensen2015#scope=Lemma2-Prop1 => Already-cited robust discount-factor comparative statics for a single agent's dynamic program and for aggregates; positioning of the one-government patience comparison, not support for a result.
% support_status: VERIFIED
% status: PROPOSED
% @end-econ-contract
\central{The threshold balances the target's net cost of sanctions against the value of
the entrenchment they build.} In words, an equilibrium with full non-compliance exists
exactly when the target's discount factor $\beta_T$ is at least $\beta_T^*$. It rises with this net cost and falls with the target's
valuation of the political state. The non-existence result applies across all
equilibria in $\mathcal E$: below $\beta_T^*$, no equilibrium exhibits full
non-compliance. When the net cost equals or exceeds the target's valuation of the
political state, $\beta_T^*\ge1$ and no discount factor can support full
non-compliance. The sender's patience changes the level of its response, and thus the
compliance threshold for relief, but in the special case it does not shift the target's
threshold, which depends on the target's costs, pain, state valuation, and discount
factor. \citet[Lemma~2, Proposition~1]{AcemogluJensen2015} derive robust
discount-factor comparative statics for single-agent dynamic programs and for
aggregates of price-taking agents; our result concerns one government's patience in a
two-player strategic setting.

% @econ-contract paragraph
% id: RESULTS_PATIENCE
% role: State that existence of a full-defiance equilibrium is inherited by a more patient Target, at its exact hypotheses, with the slope condition in words.
% depends_on: RESULTS_THRESHOLD_READING
% handoff: Why patience preserves defiance.
% exact_ref: tex-label:prop:patience-existence
% support_status: VERIFIED
% status: PROPOSED
% @end-econ-contract
\central{Outside the special case, the sufficient conditions for relief and for
non-compliance need not coincide at a single discount factor, but a more patient target
still preserves the full-non-compliance equilibrium.} The target's \emph{slope
condition} at a non-negative slope is the condition under which its value under full
non-compliance rises with the state at least at that slope. It requires that the
target's per-period stake under sanctions cover the slope effect---net of state
persistence---plus the portion of sanction pain left uncompensated by entrenchment at
that slope, scaled by the responsiveness of the sender's reaction to changes in the
state. Under the entrenchment condition, the slope condition holds at the slope that
the entrenchment condition uses.

\begin{proposition}[Greater target patience preserves the full-non-compliance equilibrium]
\label{prop:patience-existence}
Consider two economies in the class that differ only in the target's discount factor,
with $\beta_T\le\beta_T'$. Suppose the sender's response at full non-compliance at its
constructed value is non-decreasing in the state, and the slope condition holds at
some non-negative slope for the less patient economy. Then:
\begin{enumerate}[label=(\roman*)]
\item the slope condition holds at the same slope for the more patient economy; and
\item if an equilibrium in $\mathcal E$ with full non-compliance at every state exists
at $\beta_T$, one exists at $\beta_T'$.
\end{enumerate}
\end{proposition}

% @econ-contract paragraph
% id: RESULTS_PATIENCE_READING
% role: Read prop:patience-existence: why greater patience preserves defiance, and what the slope condition is for (restored from the pre-exposition text after the claim check).
% depends_on: RESULTS_PATIENCE
% handoff: The general sufficient conditions and branch invariance.
% exact_ref: tex-label:prop:patience-existence
% support_status: VERIFIED
% status: PROPOSED
% @end-econ-contract
\central{Under these two conditions, if full non-compliance is an equilibrium at a given
target discount factor, it remains an equilibrium at every higher discount factor.}
This result extends to discount factors for which neither sufficient condition holds.
The sender's value function and best response at full non-compliance are identical
across both economies. Greater target patience increases the value of the entrenchment
that compliance would forgo, since compliance lowers the sender's response; the net gain
from non-compliance is therefore preserved. The slope condition ensures that the
target's value under full non-compliance is non-decreasing in the state, as required for
the argument; assuming this monotonicity directly at either economy also suffices. The
result is a comparative statics result for the viability of a single profile---full
non-compliance at all states---as an equilibrium.

% @econ-contract paragraph
% id: RESULTS_BRACKET_WORDS
% role: State in words the general sufficient conditions for relief and defiance, the gap between them, the exact thresholds and the continuum example, and branch invariance (Tier 3 and 4; statements in the Online Appendix).
% depends_on: RESULTS_PATIENCE_READING
% handoff: The Target's minmax payoff and uniqueness.
% merged_from: RESULTS_CERTIFICATES; RESULTS_BRANCH_INVARIANCE
% exact_ref: tex-label:prop:bracket
% exact_ref: tex-label:prop:exact-thresholds
% exact_ref: tex-label:prop:affine-window
% exact_ref: tex-label:prop:branch-invariance
% support_status: VERIFIED
% status: PROPOSED
% @end-econ-contract
\central{Across the full class of economies, sufficient conditions on the target's
discount factor for relief and for non-compliance, their exact thresholds, and the
result that patience---like every other target parameter---shapes whether a
full-non-compliance equilibrium exists but not its equilibrium trajectory are all
stated in \Cref{sec:moved-main-2026-10-13:bracket} of the Online Appendix.}

\label{sec:results:security}
\label{sec:results:identification}

% @econ-contract paragraph
% id: RESULTS_SECURITY
% role: State in words the minmax bound and the uniqueness among non-decreasing sanction rules under the strong entrenchment condition (Tier 3: prop:security, prop:identification, stated in the Online Appendix), and why.
% depends_on: RESULTS_BRACKET_WORDS
% handoff: Welfare.
% merged_from: RESULTS_SECURITY_ACCOUNTING; RESULTS_IDENTIFICATION
% note: 2026-10-16 (structure plan S4): no longer its own subsection; the comparison with history-dependent strategies and the conflict literature moved here from the Discussion.
% exact_ref: tex-label:prop:security
% exact_ref: tex-label:prop:identification
% exact_ref: bib:PhelanStacchetti2001#scope=paper
% source_role: bib:PhelanStacchetti2001#scope=paper => Already-cited sequential equilibria supported by history-dependent strategies with public sunspots; the class of strategies the minmax bound covers; positioning only.
% exact_ref: bib:Powell2004#scope=paper
% source_role: bib:Powell2004#scope=paper => Already-cited conflict model stating its results through minmax payoffs over all equilibria; a comparison point, not support.
% exact_ref: bib:KraininWiseman2016#scope=paper
% source_role: bib:KraininWiseman2016#scope=paper => Already-cited all-equilibrium patience result in a conflict model at a common discount factor; a comparison point, not support for this paper's bound.
% support_status: VERIFIED
% status: PROPOSED
% @end-econ-contract
\central{Under the strong entrenchment condition, non-compliance also bounds how far
the sender can lower the target's payoff and, across the class, pins down every
equilibrium with a non-decreasing sanction rule.}\footnote{Both results are stated in
\Cref{sec:moved-main-2026-10-13} and proved in \Cref{app:proofs:scope} of the
Online Appendix.} The \emph{strong entrenchment
condition} strengthens \Cref{ass:loss-offset} by removing the
offset from entrenchment's softening of the pain of additional sanctions, making it a
more demanding requirement; only the two results in this paragraph rely on it.
Permanent relief---with the target maintaining non-compliance throughout---then
constitutes the harshest treatment the sender can impose: the resulting payoff equals
the target's minmax payoff in pure strategies. The target can guarantee itself this
payoff by maintaining non-compliance in every period against any history-dependent
sanction strategy, and every equilibrium in $\mathcal E$ yields the target at least this
payoff. History-dependent strategies of this form, together with a random signal that all
players observe, support the sequential equilibria of \citet{PhelanStacchetti2001}, and the minmax payoff is the
standard object through which the conflict literature characterizes all-equilibrium
results \citep{Powell2004,KraininWiseman2016}. The intuition is that, under the strong
form, the entrenchment built by sanctions is worth at least as much to the target as
the associated pain. Under the same condition, across the class, every equilibrium with
a sanction rule non-decreasing in the state exhibits full non-compliance, along with the
constructed values and response; no equilibrium exhibits zero sanctions at every state.
The next section turns from the governments' incentives to welfare.
